\documentclass[10pt]{article}

\usepackage[utf8]{inputenc}

\usepackage[T1]{fontenc}

\usepackage{amsmath}

\usepackage{amsfonts}

\usepackage{amssymb}

\usepackage[version=4]{mhchem}

\usepackage{stmaryrd}



\DeclareUnicodeCharacter{0131}{$\imath$}



\begin{document}

Теорема была получена С. А. Чаплыгиным в 1919. Лum.: [1] М а м е д о в Я. Д., Аш и р о в С., А т д а е в С., Теоремы о неравенствах, Аш., 1980.



См. также Јит. при статье Дияберенциальное неравенство.



ЧАСТИЦ МЕТОД - метод численного әксперимента для моделирования движения сплошных или дискретных сред. Многие Ч. м. используют эйлерово-лагранжево или лаг ранжево ошисание движущейся среды. Для решения системы уравнений сжимаемой среды напболее распространен крупных частиц метод (см. [1]), применяемый при исследовании одно- и многофазных гомои гөтерогенных потоков газа и плазмы. If Ч. м. относится мет о д с о бодны х точе к (cм. [1], [2]), в к-ром отсутствует фиксированный шаблон. Одним из первых несовершенных Ч. м. является м е т о д ч а ст и ц в я ч е ї к а х (метод РIC, см. [3]). В нем используются две расчетные сетки: эйлерова и лагранжева. Из-за дискретности представления сплошной среды этому методу присущи значительные флуктуации решений. Для уменьшения флуктуаций используется методика частиц-слоев в пространственно-одномерном случае. Методу PIC близок метод FLIC (см. [4]) обладающий плохими диссипативными свойствами При расчете несжимаемых сред используются метод МӒС (см. [5]) и метод SMAC (см. [6]), где частицы играют роль маркеров для выделения поверхностей раздела сред. Ч. м. получили распространение при описании турбулентности, в динамике разреженных газов, при решении задач әлектродинамики и др. (см. [1], [7]).



Jum.: [1] Б е л оц е р к о в с к и й O. М., Д а в ыд о в Ю. М., Метод крупных частиц в газовой динамике, М., физ.», 1965, т. 5, № 4, c. 680-88; [3] E v a n s M. W., फn3."), 1965, T. 5, № 4 , c. $680-88$; L3] E v a n s M. W., dynamic calculations, Los Alamos, 1957; [4] G e n t r y R. A., M a r t i n R. E., D a l y B. J., "J. Comp. Phys.", 1966, v. 1 ; $\mathcal{N}_{2} 1$, p. 87-118; [5] The MAC-method, Los Alamos, 1966; [6] numerical technique for calculating incompessible fluid flows, fos Alamos, 1970; [7] Б е р е з и н Ю. А., В ш̈ и в к о в В. A., Метод частиц в динамике разреженной плазмы, Новосиб.,

1980.



УАСТИЧНАЯ ГЕОМЕТРИЯ - инцидентностная структура $S=(P, L, I)$, в к-рой отногение инцидентности между точками и прямыми симметрично и удовлетворяет следующим аксиомам:



\begin{enumerate}

  \item каждая точка инцидентна $r$ прямым, $r \geqslant 2$, и две различные точки инцидентны не более чем одной прямой;



  \item каждая прямая инцидентна $k$ точкам, $k \geqslant 2$;



  \item через каждую точку, не инцидентную прямой $l$, проходит ровно $t \geqslant 1$ прлмых, пересекающих $l$.



\end{enumerate}



Если Ч. г. состоит из $v$ точек и $b$ грямых, то



$v=k[(k-1)(r-1)+t] / t \quad$ и $\quad b=r[(k-1)(r-1)+t] / t$, а необходимыми условиями существования такої Ч. г. являются делимость $(k-1)(r-1) k r$ на $t(k+r-t-1)$, $k(k-1)(r-1)$ на $t$ и $r(k-1)(r-1)$ на $t$ (см. [2]).



Ч. г. можно разделить на четыре класса:



a) Ч. г. с $t=k$ или (двойственно) $t=r$. Геометрии этого типа - то же самое, что $2-(v, k, 1)$-схема или $2-(v, r$, 1)-схема (см. Блок-схема);



б) Ч. г. с $t=k-1$ или (двойственно) $t=r-1$. В әтом случае Ч. г.- то же самое, что и сеть порядка $k$ и дефекта 1 или $k-r+1$



в) Ч. г. с $t=1$, наз. о 6 о рех угол вник а ми;



г) Ч. г. с $1<t<\min (k-1, r-1)$. Все известные Ч. г. этого типа строятся с помощью максимальных дуг в шроективных плоскостях (см., напр., [3]).



Jum.: [1] B o s e R. C., "Pacific J. Math.», 1963, v. 13, № 2, p. 389-419; [2] T a c Д̆ ж. A., "Математика", 1980, 'т. 19, c. $82-100 ;$ [3] Thas J. A., "Geometriae ded icata", 1974, v. 3,

№ 1, p. $61-64$.

B. B. Agанасвев.



ЧАСТИЧНАЯ ПРОБЈЕМА с о б с в е Н Н зн а ч е и й - задача вычисления одного или не- скольких собственных значений квадратной матрицы, обычно действительной или комшлексной, а также соответствующих им собственных векторов.



Чаще всего в практике встречаются следующие варианты Ч. п. собственных значений: 1) найти группу наименьших (наибольших) по абсолютной величине собственных значениї; 2) найти группу собственных значений, ближайших к заданному числу $a ; 3)$ найти тотки спектра, принадлежащие заданному интервалу $(\alpha, \beta)$ (для симметричной или әрмитовой матрицы).



Большинство методов решения Ч. П. собственных значений для матриц общего вида имеет в основе идею степенной итерации или ее разновидность- обратную итерацию (см. Итерационные методы решения проблемы собственных значений матрицы): если матрица $A$ обладает доминирующим по абсолютной величине собственным значением $\lambda_{\max }$ и $x_{\max }$ - соответствующий собственный вектор, то для почти любого вектора $v$ последовательность $v, A v, A^{2} v, \ldots, A^{k} v, \ldots$ сходится по направлению к вектору $x_{\max }$. Если требуется наименьшее по абсолютной величине собственное значение (задача 2)), то степенная итерация проводится с матрицей $A^{-1}$ (метод обратной итерации); при вычислении собственного значения, ближайшего к $a$ (задача 3)),с матрицей $(A-a)^{-1}$ (метод обратной итерации со сдвигом).



Наиболее ванкый частный случай Ч. п. собственных значений - вычисление собственных значений и собственных вөкторов действительной симметричной либо комплексной әрмитовой матрицы $A$. Здесь имеется ряд әффективных численных методов решения Ч. П. собственных значений, основанных на весьма различных идеях (см. [1]). Среди них: методы, использующие экстремальные свойства функционала Рэлея (наибольшее и наименьшее собственные значения матрицы $A$ реализуют соответственно максимум и минимум отношения Рәлея $\varphi(A, z)=(A z, z) /(z, z), \quad z \neq 0 ; \quad$ достигаются эти экстремумы на соответствующих собственных векторах); методы, исгользующие закон инерции Сильвестра (метод юоследовательности Штурма и, более общо, методы деления спектра); наконец, методы, базирующиеся на аıпроксимационных свойствах крыловских подпространств, т. е. линейных оболочек систем вида $v, A v, \ldots, A^{k-1} v$ (метод Ланцона и его варианты). Выбор того или иното метода определяется такими соображениями, как порядок задачи, стенень разреженности матрицы, наличие ленточноӥ структуры, доступная информация о спектро и т. Д.



Методы репения Ч. ІІ. собственных значений как в общем, так и в эрмштовом случае можно разделить на групповые и последовательные. Групповые методы характеризуются тем, что в них нужные собственные значения (и соответствующие собственные векторы) вычисляются в известной мере параллельно. Сюда относятся многочисленные методы одновременных итераций, метод Јанцоша, методы деления спектра.



В последовательных методах собственные значения определяются поочередно. При этом, начиная со второго собственного знатения, возникае' необходимость воспрепятствовать тому, Ітобы итерации сходились к ранее найденным корням. С этоїі необходимостью связаны различные приемы исчерпывания (или дефляции) [2]. В одних случаях исqерпывание приводит к построөнию матрицы $\tilde{A}$, у к-роӥ вычисленным собственным значениям $A$ соответствуют нулевыс корни; в остальном спектр обеих матриц совнадает, совгадают и их собственные векторы. В других случаях результатом исчерпывания является расцепление матрицы, вследствие чего последовательные собственные значения можно ощределить, пользуясь матрицами убывающих порядков. В третьих - итерации метода с неизменної матрицеӥ $\boldsymbol{A}$ сопровогдаотся ортогонализацией по отноше-





\end{document}